\documentclass[11pt,letterpaper]{article}

% =====================================================================
% Homework 1 - Theory part
% Course: Deep Learning - Universidad del Pacifico (UP)
% Instructor: Alexander Quispe
% =====================================================================

\usepackage[utf8]{inputenc}
\usepackage[T1]{fontenc}
\usepackage[english]{babel}
\usepackage{lmodern}
\usepackage[margin=2.5cm]{geometry}
\usepackage{amsmath, amssymb, amsthm, mathtools}
\usepackage{enumitem}
\usepackage{graphicx}
\usepackage{xcolor}
\usepackage{tikz}
\usetikzlibrary{positioning, arrows.meta, calc}
\usepackage[hidelinks]{hyperref}
\usepackage{fancyhdr}
\usepackage{titlesec}

% --- Style ----------------------------------------------------------
\pagestyle{fancy}
\fancyhf{}
\fancyhead[L]{\small Deep Learning -- UP}
\fancyhead[R]{\small Homework 1 (Theory)}
\fancyfoot[C]{\thepage}
\renewcommand{\headrulewidth}{0.4pt}

\titleformat{\section}{\large\bfseries}{Problem \thesection.}{0.5em}{}
\setlength{\parindent}{0pt}
\setlength{\parskip}{0.4em}

% --- Macros ---------------------------------------------------------
\newcommand{\R}{\mathbb{R}}
\newcommand{\E}{\mathbb{E}}
\newcommand{\Var}{\mathrm{Var}}
\newcommand{\sgn}{\mathrm{sgn}}
\newcommand{\softmax}{\mathrm{softmax}}
\newcommand{\relu}{\mathrm{ReLU}}
\newcommand{\diag}{\mathrm{diag}}
\DeclareMathOperator*{\argmin}{arg\,min}

\newcommand{\pts}[1]{\textit{(#1 pts)}\;}
\newcommand{\note}[1]{\textit{\small Note: #1}}

% =====================================================================
\title{\vspace{-1.5cm}\textbf{Homework 1 -- Theory Part}\\[0.3em]
       \large Deep Learning -- Universidad del Pac\'ifico}
\author{Alexander Quispe}
\date{Due date: TBD}

\begin{document}
\maketitle
\thispagestyle{fancy}

\section*{Instructions}

\begin{itemize}[leftmargin=1.5em]
  \item This part covers topics from \textbf{Lectures 1, 2 and 3}, complemented
        with exercises adapted from the Caltech \emph{Learning From Data}
        Problem Set 4 (Yaser S. Abu-Mostafa).
  \item There are \textbf{6 problems} worth \textbf{100 points} in total.
  \item Justify every step. A correct final answer without derivation does not
        receive full credit.
  \item You may submit handwritten (clear scan/photo) or typeset in LaTeX.
        LaTeX is preferred.
  \item For network-design problems (3 and 4) you may include a diagram
        (TikZ, hand-drawn scan, or equivalent).
\end{itemize}

\begin{center}
\begin{tabular}{|c|l|c|}
\hline
\textbf{\#} & \textbf{Topic} & \textbf{Points} \\
\hline
1 & Manual backpropagation of a 2-layer network        & 21 \\
2 & Activation functions and the vanishing gradient    & 16 \\
3 & Building OR with ReLU (Caltech)                    & 12 \\
4 & Building XOR with ReLU (Caltech)                   & 19 \\
5 & Weight initialization (Xavier and He)              & 16 \\
6 & The Adam optimizer: bias correction                & 16 \\
\hline
\multicolumn{2}{|r|}{\textbf{Total}} & \textbf{100} \\
\hline
\end{tabular}
\end{center}

\hrulefill

% =====================================================================
\section{Manual backpropagation of a 2-layer network \pts{21}}

Consider a fully connected network with the following architecture:

\begin{itemize}[leftmargin=1.5em, itemsep=0pt]
  \item Input: $\mathbf{x} \in \R^{d}$.
  \item Hidden layer: $\mathbf{z}^{(1)} = W^{(1)} \mathbf{x} + \mathbf{b}^{(1)}$,
        with $W^{(1)} \in \R^{H \times d}$, $\mathbf{b}^{(1)} \in \R^{H}$, and
        activation $\mathbf{h} = \sigma(\mathbf{z}^{(1)})$ where $\sigma$ is the
        sigmoid applied element-wise.
  \item Output layer: $z^{(2)} = \mathbf{w}^{(2)\top} \mathbf{h} + b^{(2)}$,
        with $\mathbf{w}^{(2)} \in \R^{H}$, $b^{(2)} \in \R$, and prediction
        $\hat{y} = \sigma(z^{(2)})$.
  \item Loss: \emph{binary cross-entropy}
        $\mathcal{L} = -\,y \log \hat{y} - (1-y)\log(1-\hat{y})$, with $y \in \{0,1\}$.
\end{itemize}

\textbf{(a)} \pts{5} Show that for this combination (sigmoid + BCE):
$$
\frac{\partial \mathcal{L}}{\partial z^{(2)}} \;=\; \hat{y} - y.
$$

\textbf{(b)} \pts{7} Compute the gradients
$\dfrac{\partial \mathcal{L}}{\partial \mathbf{w}^{(2)}}$ and
$\dfrac{\partial \mathcal{L}}{\partial b^{(2)}}$.
State the shape of each tensor.

\textbf{(c)} \pts{7} Compute the gradients
$\dfrac{\partial \mathcal{L}}{\partial W^{(1)}}$ and
$\dfrac{\partial \mathcal{L}}{\partial \mathbf{b}^{(1)}}$ via the chain rule.
Identify the factor $\sigma'(\mathbf{z}^{(1)})$ and express it as
$\mathbf{h} \odot (1 - \mathbf{h})$.

\textbf{(d)} \pts{2} Argue why \emph{a single} backward pass is enough to
compute every gradient, and why a finite-difference estimate would be
$\mathcal{O}(P)$ times more expensive, where $P$ is the total number of
parameters.

% =====================================================================
\section{Activation functions and the vanishing gradient \pts{16}}

\textbf{(a)} \pts{5} Compute the derivatives of the following activations,
expressing the result \emph{in terms of the activation itself} whenever possible:
\begin{itemize}[leftmargin=1.5em, itemsep=0pt]
  \item $\sigma(z) = \dfrac{1}{1+e^{-z}}$
  \item $\tanh(z)$
  \item $\relu(z) = \max(0, z)$
\end{itemize}

\textbf{(b)} \pts{3} Show that $\sigma'(z)$ attains its maximum at $z=0$ with
value $1/4$. Conclude that in a deep \textbf{sigmoid} network the magnitude of
the gradient reaching the first layer is upper bounded by
$$
\left\| \nabla_{W^{(1)}} \mathcal{L} \right\| \;\leq\; C \cdot (1/4)^{L},
$$
for some $C > 0$ that depends on the remaining weights, where $L$ is the number
of hidden layers. Briefly discuss the practical implication.

\textbf{(c)} \pts{3} For $\tanh$ the analogous factor is $1$ (at $z=0$).
\textbf{Why}, in practice, does $\tanh$ still fail to fully solve the vanishing
gradient problem in very deep networks?

\textbf{(d)} \pts{5} ReLU avoids the vanishing gradient but introduces the
\emph{Dying ReLU} problem.
\begin{enumerate}[label=(\roman*), leftmargin=2em, itemsep=2pt]
  \item Formally define when a ReLU unit is \emph{dead}.
  \item If a unit is dead over the entire dataset, \textbf{what is its
        gradient} with respect to $W^{(1)}$? What does this imply for the
        weight update?
  \item Mention \emph{two} standard strategies (from class or from the Caltech
        problem set 4) to mitigate this problem.
\end{enumerate}

% =====================================================================
\section{Building OR with ReLU \pts{12}}

\note{Adapted from Caltech Problem Set 4, Section 1.D.}

Design a \textbf{fully connected ReLU network} that implements the OR function
on binary inputs $x_1, x_2 \in \{0, 1\}$. The output $f(x_1, x_2)$ must satisfy:
$$
\begin{aligned}
\text{OR}(0,0) &= 0, & \text{OR}(0,1) \;\geq\; 1, \\
\text{OR}(1,0) &\geq 1, & \text{OR}(1,1) \;\geq\; 1.
\end{aligned}
$$

\textbf{(a)} \pts{7} Provide the exact weights and biases. Use the
\textbf{minimum number of hidden units} possible (justify the minimum). You may
draw the network or write the matrices/vectors.

\textbf{(b)} \pts{5} Verify your network on the truth table by explicitly
evaluating it on the four inputs and showing the per-layer computation.

% =====================================================================
\section{Building XOR with ReLU \pts{19}}

\note{Adapted from Caltech Problem Set 4, Section 1.E.}

Design a \textbf{fully connected ReLU network} that implements the XOR
function on binary inputs:
$$
\begin{aligned}
\text{XOR}(0,0) &= 0, & \text{XOR}(0,1) \;\geq\; 1, \\
\text{XOR}(1,0) &\geq 1, & \text{XOR}(1,1) \;=\; 0.
\end{aligned}
$$

\textbf{(a)} \pts{5} Prove formally that \textbf{no} network \emph{without
hidden layers} (i.e., a single linear layer with a ReLU at the end) can
implement XOR.
\note{Hint: a function of the form $\relu(\mathbf{w}^\top \mathbf{x} + b)$ is
convex in $\mathbf{x}$. \textbf{Is XOR convex?}}

\textbf{(b)} \pts{5} State the \textbf{minimum number of hidden layers} that
allows a ReLU network to implement XOR, and justify.

\textbf{(c)} \pts{7} Build an explicit minimal network: specify weights,
biases, and architecture.

\textbf{(d)} \pts{2} Verify your construction on the four inputs by showing
the computation.

% =====================================================================
\section{Weight initialization: Xavier and He \pts{16}}

Consider a linear layer $\mathbf{z} = W \mathbf{x}$ with
$W \in \R^{n_{\text{out}} \times n_{\text{in}}}$, and assume the inputs $x_i$
are i.i.d.~with zero mean and variance $\sigma_x^2$, and the weights $W_{ij}$
are i.i.d.~with zero mean, variance $\sigma_W^2$, and \textbf{independent of
$\mathbf{x}$}.

\textbf{(a)} \pts{5} Show that the variance of each component $z_k$ satisfies
$$
\Var(z_k) \;=\; n_{\text{in}}\, \sigma_W^2\, \sigma_x^2.
$$
\note{Recall that $\Var(\sum_i A_i) = \sum_i \Var(A_i)$ when the $A_i$ are independent.}

\textbf{(b)} \pts{3} For $\Var(z_k) = \Var(x_i)$ to hold (i.e., for the signal
to neither blow up nor vanish through the layer) it suffices to set
$\sigma_W^2 = 1/n_{\text{in}}$. If we additionally want to preserve variance
\emph{in the backward pass}, we obtain $\sigma_W^2 = 1/n_{\text{out}}$. The
\textbf{Xavier/Glorot} initialization averages both conditions:
$$
\sigma_W^2 \;=\; \frac{2}{n_{\text{in}} + n_{\text{out}}}.
$$
In 3-5 lines, discuss \textbf{why} this compromise makes sense and which
activation functions are its \emph{ideal case}.

\textbf{(c)} \pts{5} For \textbf{ReLU} the derivation changes. Assuming
$\mathbf{x}$ comes from a previous ReLU layer (so on average half of the
components are zero) and that the pre-activation is symmetric around zero,
show that preserving variance requires
$$
\sigma_W^2 \;=\; \frac{2}{n_{\text{in}}}\quad \text{(He et al., 2015)}.
$$
\note{Key fact: $\E[\relu(z)^2] = \tfrac{1}{2}\,\Var(z)$ when $z \sim \mathcal{N}(0, \Var(z))$.}

\textbf{(d)} \pts{3} If you used \textbf{Xavier instead of He} in a deep ReLU
network, what would happen to the variance of the activations as depth
increases? Reason by computing the per-layer compression factor.

% =====================================================================
\section{The Adam optimizer: bias correction \pts{16}}

Adam combines momentum and RMSProp with a \emph{bias correction} step. The
updates are:
\begin{align*}
\mathbf{m}_t &= \beta_1\, \mathbf{m}_{t-1} + (1-\beta_1)\, \mathbf{g}_t, \\
\mathbf{v}_t &= \beta_2\, \mathbf{v}_{t-1} + (1-\beta_2)\, \mathbf{g}_t^{\,2}, \\
\hat{\mathbf{m}}_t &= \frac{\mathbf{m}_t}{1 - \beta_1^{\,t}}, \qquad
\hat{\mathbf{v}}_t = \frac{\mathbf{v}_t}{1 - \beta_2^{\,t}}, \\
\boldsymbol{\theta}_t &= \boldsymbol{\theta}_{t-1}
        - \alpha \cdot \frac{\hat{\mathbf{m}}_t}{\sqrt{\hat{\mathbf{v}}_t} + \varepsilon},
\end{align*}
with $\mathbf{m}_0 = \mathbf{v}_0 = \mathbf{0}$ and $\beta_1, \beta_2 \in [0, 1)$.

\textbf{(a)} \pts{5} Assume gradients are stationary, $\E[\mathbf{g}_t] = \mathbf{g}$
for every $t$. Prove by induction that
$$
\E[\mathbf{m}_t] \;=\; (1 - \beta_1^{\,t})\, \mathbf{g}.
$$
Conclude that \textbf{without the correction} $\hat{\mathbf{m}}_t = \mathbf{m}_t / (1 - \beta_1^t)$
the first iterations are biased toward zero.

\textbf{(b)} \pts{5} Analogously, assuming $\E[\mathbf{g}_t^{\,2}] = \mathbf{v}$
for every $t$, prove that
$$
\E[\mathbf{v}_t] \;=\; (1 - \beta_2^{\,t})\, \mathbf{v}.
$$

\textbf{(c)} \pts{3} With $\beta_1 = 0.9$, $\beta_2 = 0.999$, numerically
evaluate the factors $1 - \beta_1^{\,t}$ and $1 - \beta_2^{\,t}$ for
$t \in \{1, 5, 50, 500\}$ and comment on the regime where the correction
\emph{stops mattering}.

\textbf{(d)} \pts{3} Compare Adam with \textbf{SGD + Momentum}. Give
\emph{one} situation where SGD+Momentum is preferable, and \emph{one} where
Adam is. Justify both in terms of loss-surface geometry or data structure.

% =====================================================================
\vspace{1em}
\hrulefill

\begin{center}
\textit{End of Homework 1 -- Theory Part}
\end{center}

\end{document}
